\section{Introduction}
Knowledge representation in structured form has become an increasingly important pillar in the development of intelligent systems. A key paradigm is the Knowledge Graph (KG), which represents real-world entities as nodes and the relationships between them as edges, enabling both human-readable and machine-processable semantics. Parallel to this, the advent of large language models (LLMs) has transformed how unstructured textual data is used and processed. However, while LLMs excel at language understanding, they are prone to producing content that is not grounded in factual evidence or reliable structure, commonly termed \textit{hallucinations}. Integrating knowledge graphs with LLMs can help ground their outputs in structured knowledge and reduce unanchored generation.
Documents in PDF format remain one of the most common mediums in many domains (academic, technical, regulatory). Yet, extracting structured knowledge from PDFs poses substantial technical challenges: the text may be unstructured, include tables and figures, hierarchical headings, and embedded images.

In this context, our project addresses the following research problem: how to extract structured knowledge from a domain-specific PDF (in the food domain), construct a knowledge graph (KG), and use that graph to generate grounded instruction-response pairs in order to fine-tune an LLM, and then compare this fine-tuned model with two baselines: (i) zero-shot LLM, and (ii) retrieval-augmented generation over the KG (RAG-KG). 

Recent dietary guidelines such as the Finnish “National Nutrition Recommendations 2024” emphasise not only individual health, but also the planetary health consequences of food production and consumption. For instance, food consumption in Finland has been estimated to account for around $20 \%$ of the annual climate impacts of the average consumer, and diets aligned with nutritional and sustainability goals can potentially reduce the climate footprint by over one-third. These guidelines are produced through large-scale surveys, modelling of consumption patterns, and systematic reviews of nutritional and environmental research.
As the main source of domain knowledge, this project relies on the “Sustainable Health from Food – National Nutrition Recommendations 2024” report by the Finnish Institute for Health and Welfare \cite{thl2024nutrition}. This document provides an evidence-based synthesis of food and health recommendations, yet it exists primarily as a human-readable PDF, where valuable knowledge is distributed across sections, tables, and figures.
In this project, we take this complex PDF as our source document and convert its contents into structured knowledge within the food and nutrition domain. By doing so, we establish a pipeline that enables downstream tasks—such as fact-based instruction–response generation and fine-tuning of large language models—and allow for systematic comparison between grounded and ungrounded model settings.


